\section{18.09.2026} 

\subsection{Задачи}

\subsubsection{Задача 3}

\begin{exercise}
    Из последовательности чисел $1, 2, \ldots, n$ наудачу выбирается 2 числа. Какова вероятность, что одно из них меньше $k$, а другое больше $k$, где $1 < k < n$ --- произвольное целое число.
\end{exercise}

\begin{solution}
    \[
        \Pr(\mathcal{A}) = \frac{N_{\text{благоприятных}} }{N_{\text{общих}}} = \frac{(k-1)(n-k)}{C^2_n}
    \]
\end{solution}

\subsubsection{Задача 4}

\begin{exercise}
    $n$ друзей садятся случайным образом за круглый стол. Найти вероятность того, что
    \begin{itemize}
        \item[а)] два фиксированных лица $A$ и $B$ сядут рядом, причем $B$ слева от $A$;
        \item[б)] три фиксированных лица $A, B$ и $C$ сядут рядом, причем $A$ справа от $B$, а $C$ --- слева 
    \end{itemize}
\end{exercise}

\begin{solution}
    В таких задача важно сначала зафиксировать одного человека, и только после начинать решение.

\[
    \Pr(\mathcal{B}) = \frac{N_{\text{благоприятных}} }{N_{\text{общих}}} = \frac{(n-2)!}{(n-1)!} = \frac{1}{n - 1}
\]

--- кол-во зафиксированных по двое это кол-во благоприятных

\[
    \Pr(\mathcal{C}) = \frac{N_{\text{благоприятных}} }{N_{\text{общих}}} = \frac{(n-3)!}{(n-1)!} = \frac{1}{(n - 1)(n - 2)}
\]
\end{solution}

\subsubsection{Задача 5}

\begin{exercise}
    В урне находится 10 белых и 6 черных шаров. Без возвращения вытягиваются 4 шара. Найти вероятность того, что среди них будет ровно два черных шара; хотя бы один черный шар.
\end{exercise}

\begin{solution}
    1) Найдем вероятность того, что среди вытянутых шаров будет ровно 2 черных (событие $A$). В этом случае остальные 2 шара должны быть белыми. 
    
    Общее число возможных исходов при вытягивании 4 шаров из 16 (10 белых + 6 черных) равно числу сочетаний:
    \[
        N_{\text{общих}} = C_{16}^4 = \frac{16!}{4!(16-4)!} = \frac{16 \cdot 15 \cdot 14 \cdot 13}{1 \cdot 2 \cdot 3 \cdot 4} = 1820
    \]
    Число благоприятных исходов находим как произведение числа способов выбрать 2 черных шара из 6 и 2 белых из 10:
    \[
        N_{\text{благоприятных}} = C_6^2 \cdot C_{10}^2 = \frac{6 \cdot 5}{1 \cdot 2} \cdot \frac{10 \cdot 9}{1 \cdot 2} = 15 \cdot 45 = 675
    \]
    Искомая вероятность:
    \[
        \Pr(A) = \frac{C_6^2 \cdot C_{10}^2}{C_{16}^4} = \frac{675}{1820} = \frac{135}{364} \approx 0.371
    \]

    2) Найдем вероятность того, что вытянут хотя бы один черный шар (событие $B$). 
    
    Удобнее перейти к противоположному событию $\overline{B}$ — среди вытянутых шаров нет ни одного черного, то есть все 4 шара оказались белыми. 
    
    Число исходов, благоприятствующих событию $\overline{B}$:
    \[
        N_{\overline{B}} = C_{10}^4 = \frac{10 \cdot 9 \cdot 8 \cdot 7}{1 \cdot 2 \cdot 3 \cdot 4} = 210
    \]
    Вероятность противоположного события:
    \[
        \Pr(\overline{B}) = \frac{C_{10}^4}{C_{16}^4} = \frac{210}{1820} = \frac{21}{182} = \frac{3}{26}
    \]
    Тогда вероятность того, что попадется хотя бы один черный шар:
    \[
        \Pr(B) = 1 - \Pr(\overline{B}) = 1 - \frac{3}{26} = \frac{23}{26} \approx 0.885
    \]
\end{solution}

\subsubsection{Задача 6*}

\begin{exercise}
    Колода игральных карт содержит 52 карты, разделяющиеся на 4 различные масти по 13 карт в каждой. Предположим что колода тщательно стасована, так что вытаскивание любой карты одинаково вероятно. Вытащим 6 из них. 

\begin{itemize}
    \item[(a)] Найти вероятность того, что среди этих карт будет король пик. 
    \item[(b)] Найти вероятность того, что среди этих карт будут представители всех мастей.
\end{itemize}
\end{exercise}

\begin{solution}
    \[
        \frac{C^5_{51}}{C^{6}_{52}} = \frac{\frac{51!}{5! \cdot 46!}}{\frac{52!}{6! \cdot 46!}} = \frac{51! \cdot 6! \cdot 46!}{5! \cdot 46! \cdot 52!} = \frac{6}{52} = \frac{3}{26} \approx 0.12
    \]
    \[
        \frac{4 \cdot C^3_{13} \cdot 13^3 + C^2_4 \cdot (C^2_{13})^2 \cdot 13^2}{C^6_{52}}
    \]
\end{solution}

\subsubsection{Листок 2. Задача 1}

\begin{exercise}
    Староста 258 группы заметил, что семинарист приходит на занятия со случайным опозданием в пределах 15 мин., при этом он разрешает проходить в аудиторию только тем студентам, которые пришли после него не позднее 5 мин. Староста же решил опоздать на семинар, но выбрал себе границу случайного опоздания всего в 10 мин. и время ожидания семинариста тоже 10 мин. Какова вероятность того, что он все же посетит семинар.
\end{exercise}

\begin{solution}
    Вставить картинку

    \[
        0 < x_1 < 15
    \]
    \[
        0 < x_2 < 10
    \]
    \[
        \begin{cases}
            x_2 - x_1 < 5\\
            x_1 - x_2 < 10
        \end{cases}
    \]
    \[
        x_1 - 10 < x_2 < x_1 + 5
    \]
\end{solution}
